This study involved secondary analysis of fully de-identified electronic health record data. \new{Because the analysis involves only fully de-identified secondary data that are not readily identifiable, it does not constitute human-subjects research under the U.S. Common Rule (45~CFR~46.104) and did not require Institutional Review Board approval or a formal exemption determination at the authors' institution (Harvard Medical School); no patient contact, intervention, or access to identifiable private information occurred. The datasets were obtained and used under their respective data-use agreements: PhysioNet credentialed access for MIMIC-III (CITI Human Subjects Research training); PhysioNet credentialed access for CORAL (CITI Data or Specimens Only Research training); and a project-specific Data Use Agreement with the 4CE clinical-notes custodians at the Harvard Medical School DBMI Data Portal. All annotators completed the CITI training course required for each dataset and held the access credentials required for each dataset prior to data access. All large-language-model inference was performed via Harvard Medical School Information Technology's managed Azure AI service (HMS Azure AI),\footnote{\url{https://it.hms.harvard.edu/service/azure-ai}} an HMS-administered deployment providing project-scoped access to Microsoft Azure OpenAI models for Harvard data-security Level~4 projects; only de-identified clinical text (containing no HIPAA-regulated identifiers) was processed, and no text was retained or used to train provider models.} Ethical approvals and data-use agreements for the original data collection were obtained by the respective data custodians.